\section{综合对照：Attention alternatives 与 MoE 的共同结构}
\begin{longtable}{p{0.20\textwidth}p{0.34\textwidth}p{0.36\textwidth}}
\toprule
\textbf{方向} & \textbf{省什么} & \textbf{新增什么问题} \\
\midrule
Linear/state attention & 省完整 \(n^2\) attention matrix 和长上下文读写 & 状态表达力、长程依赖、与 full attention 的混合比例。 \\
Sparse/local attention & 省远距离 token 的 attention 成本 & 全局信息如何周期性传播，哪些 token 应该被选中。 \\
MoE & 省每 token 激活全部参数的 FLOPs & routing、load balancing、all-to-all、expert undertraining、fine-tuning 稳定性。 \\
MLA/MTP 等扩展 & 省 KV cache 或改善多 token 预测 & 与 RoPE/attention 几何、推理系统、训练目标的兼容性。 \\
\bottomrule
\end{longtable}
\begin{importantbox}{最终 takeaway}
Lecture 4 的核心不是“linear attention vs MoE 谁更好”，而是学会看 selective computation 的三件事：选择规则是什么、节省了哪种资源、为了这个节省引入了什么新的系统或训练问题。
\end{importantbox}
